\section{Problem Definition and Research Direction}\label{sec:problem-direction}

In a standard Digital Audio Workstation (DAW), the primary method of input is the piano roll---a grid where notes are
drawn manually across a timeline.
This approach is highly intuitive but fundamentally destructive from a structural perspective.
If a composer wishes to transpose a repeating 16-bar melody, change its rhythm, or conditionally alter its ending on the
final repetition, they must manually select and modify hundreds of individual note events.
The logical intent (e.g., ``repeat this phrase four times but play the snare on the off-beat during the last
repetition'') is lost, replaced by a static collection of data points.

Programmers often turn to algorithmic composition to recover that structure.
However, existing domain-specific languages for music are largely tailored toward \textit{live coding} and
\textit{real-time performance}.
Systems such as Sonic Pi, TidalCycles, Strudel, and ChucK prioritize low-latency scheduling, threaded or concurrent
execution, and continuous audio synthesis~\cite{Aaron_2016,McLean_Wiggins_2010,Roos_McLean_2023,Wang_Cook_2003}.
While powerful for improvisation and on-stage manipulation, these environments introduce concerns related to thread
synchronization, execution latency, and audio-engine configuration that are orthogonal to offline, deterministic
arrangement.

Furthermore, general-purpose programming languages equipped with MIDI libraries (e.g., Python with \texttt{mido}) offer
offline generation but lack domain-specific abstractions.
The composer is burdened with calculating explicit delta-times and manually constructing low-level binary messages,
shifting the focus away from musical creativity and toward byte-level engineering.

\subsection{Determinism and MIDI Portability}\label{subsec:determinism-portability}

Motivo targets a different execution model.
The compiler runs entirely offline: given the same source program, it produces the same MIDI event stream on every
invocation.
There is no interpreter, scheduler, or synthesis server involved at playback time.
The resulting Standard MIDI File can be opened in any DAW, notation tool, or hardware sequencer that accepts Format~1
MIDI, without requiring Motivo Studio or any Motivo runtime to be present.

This determinism is deliberate.
Live-coding systems must reconcile human reaction time, audio-buffer sizes, and network latency when streaming events to
a synthesizer.
Motivo instead resolves all temporal relationships during compilation.
The composer receives a fixed artifact---a \texttt{.mid} file---whose note placements and track structure are fully
materialized before any playback occurs.
The trade-off is explicit: Motivo does not support on-the-fly code modification during performance, but it guarantees
reproducible output suitable for downstream production workflows.

\subsection{Problem Statement}\label{subsec:problem-statement}

There exists a gap in the current ecosystem of music programming tools: a lack of a domain-specific language optimized
for \textit{offline, structural composition} that produces standard, interoperable artifacts.
Specifically, the challenge lies in creating a system that allows composers to define musical ideas as reusable,
parameterized algorithms while keeping the resulting temporal structure deterministic, free of real-time
execution jitter,
and readable by many standard production tools.

\subsection{Proposed Solution}\label{subsec:proposed-solution}

This thesis proposes the design and implementation of a compiled domain-specific language
for musical
composition.
The system aims to bridge the gap between algorithmic logic and the rigid binary specifications of the MIDI protocol.

The proposed solution is organized around four design principles.

First, music is treated as a recursive hierarchy. The language abstracts musical motifs into functions (patterns) that
can be parameterized, reused, and nested, reflecting the true hierarchical nature of musical form.

Second, the source language uses a C-style imperative flow. By utilizing familiar control flow structures
(\texttt{for}, \texttt{loop},
\texttt{if}), Motivo allows composers to procedurally generate variations and repetitions without learning esoteric
functional paradigms.

Third, the compiler follows a zero-runtime compilation model. The compiler operates strictly offline.
It performs semantic validation and then uses a stateful lowering engine to completely unroll the composition's
logic into a flat timeline.
This keeps temporal resolution deterministic within the compiler pipeline without relying on an interpreter
during playback.

Fourth, MIDI is used as the semantic target. The compiler serializes the final timeline into a Standard MIDI File
(Format 1).
This supports practical interoperability, allowing the generated compositions to be imported into many DAWs
and MIDI-aware tools for
further orchestration and mixing.
